\section{Entities and Ownership}
\label{sec:ch06-entities-and-ownership}
\index{entities|(}%

An entity is a type the router can ask for by name. That is the definition that
matters, and everything else follows from it.

\index{entities!key directive@\code{@key} directive}%
A type becomes an entity when a subgraph puts \code{@key} on
it~\autocite{apollo2026entities}. The
\code{fields} argument names the fields that together make up the key, which is
how the router will identify one instance when it needs to ask another subgraph
for more information.

\begin{graphqlsdl}
type Product @key(fields: "id") {
  id: ID!
  title: String!
  sku: String!
}
\end{graphqlsdl}

That \code{Product} exists in the Catalog subgraph, and the key says the router
can reach it by \code{id}. If a query selects \code{title} and
\code{Product.reviews} in the same document, the router fetches the
\code{title} and the \code{id} from Catalog, then sends the \code{id} to Reviews
with a query that asks for the reviews of that exact product.

The subgraph that defines the key is the entity's owner, but ownership here is
not a claim about the domain. It is a claim about who is the authority for the
entity's existence. Catalog owns \code{Product} because it is where a
\code{Product} is created; Reviews contributes fields to one but does not get to
say which products exist. That distinction matters the moment two subgraphs
disagree about whether a product is real, and chapter~\ref{ch:entity-resolution-done-right}
is where that disagreement is handled.

\subsection*{Extending an entity from another subgraph}

\index{entities!extending}%
A subgraph that wants to add fields to an entity another subgraph owns declares
the type again, with the same key, and adds its fields:

\begin{graphqlsdl}
# In the Reviews subgraph
type Product @key(fields: "id") {
  id: ID!
  reviews: [Review!]!
}
\end{graphqlsdl}

The key must match the owning subgraph's key, and the key fields must appear
here too, because the router has to fill them in from the representation it
already holds. \code{id} appears in both subgraphs, in both schemas, and in the
supergraph it appears once.

\index{entities!reference resolver}%
That arrangement works because the Reviews subgraph can answer this question:
given a \code{Product} identifier, give me back a \code{Product} with the
reviews attached. The function that answers it is the reference resolver, and
the router calls it through the \code{\_entities} query when it needs to bridge
fields from two subgraphs. A reference resolver receives a list of
representations, each a JSON object with \code{\_\_typename} and the key fields,
and returns the matching entities. Chapter~\ref{ch:how-a-federated-query-actually-runs}
shows the exact HTTP traffic; for the model, what matters is that a subgraph
extending an entity promises to resolve it from its key.

\subsection*{Multiple keys on one type}

\index{entities!multiple keys}%
An entity can declare more than one key, and the second one enables something
the first cannot:

\begin{graphqlsdl}
type Product @key(fields: "id") @key(fields: "sku") {
  id: ID!
  sku: String!
  title: String!
}
\end{graphqlsdl}

Now the router has two routes to the same entity. When a query starts from a
field that only hands back a \code{sku}, the router can still reach the rest of
\code{Product} without first resolving \code{id}. The downstream cost is that
every subgraph extending the entity has to be able to resolve it from either
key, which means a reference resolver for each. Chapter~\ref{ch:entity-resolution-done-right}
covers that cost and the batching strategy that makes it acceptable.

The second key does not create a second entity. It creates a second way to reach
the same one, and the router picks based on what it already has.

\subsection*{What is not an entity}

\index{entities!value types}%
A type without \code{@key} is a value type. It exists only in the subgraph that
defines it, cannot be extended by another subgraph, and the router has no way to
ask for it by name. Money, an address, a cursor, a page of results: these are
value types, and most types in a federated graph are value types. An entity is
the exception, and it is expensive specifically because it crosses boundaries.
Mark too many things as entities and you have promised the router can ask for
them, which means every subgraph that touches them has to write a reference
resolver. Chapter~\ref{ch:hard-modeling-problems} is about where the boundary
between entity and value type should be drawn, and the short version is that
most boundary types are entities and most nested types are not.

\index{entities|)}%
